\documentclass[11pt,letterpaper]{article}
\usepackage[T1]{fontenc}
\usepackage[utf8]{inputenc}
\usepackage[margin=0.88in,top=0.84in,bottom=0.83in,headheight=13pt,headsep=20pt,footskip=28pt]{geometry}
\usepackage{newpxtext}
\usepackage{amsmath,amsthm}
\usepackage{newpxmath}
\usepackage{microtype}
\usepackage{xcolor}
\definecolor{Forest}{HTML}{30392C}
\definecolor{Olive}{HTML}{687144}
\definecolor{Muted}{HTML}{65685F}
\definecolor{Sage}{HTML}{CBD0BC}
\definecolor{Pale}{HTML}{F2F4EB}
\usepackage{mathtools}
\usepackage{booktabs,tabularx,longtable,array}
\usepackage{enumitem}
\usepackage{titlesec}
\usepackage{fancyhdr}
\usepackage[most]{tcolorbox}
\usepackage{needspace}
\PassOptionsToPackage{obeyspaces}{url}
\usepackage{xurl}
\usepackage[colorlinks=true,linkcolor=Olive,citecolor=Olive,urlcolor=Olive,bookmarksnumbered=true,pdfencoding=auto,psdextra]{hyperref}
\hypersetup{pdftitle={The measurable strength of cofinal-orbit rigidity},pdfauthor={Prepared for Vladimir Reshetnikov},pdfsubject={Equiconsistency with a measurable cardinal and canonical forcing decompositions},pdfkeywords={rigid cofinal classes, measurable cardinal, Prikry forcing, HOD, Vopenka algebra, Dodd Jensen covering}}
\urlstyle{same}
\setlength{\parindent}{0pt}
\setlength{\parskip}{5.1pt plus 1pt minus .4pt}
\linespread{1.035}
\setlength{\emergencystretch}{2em}
\setcounter{tocdepth}{1}
\setcounter{secnumdepth}{2}
\titleformat{\section}{\Large\bfseries\color{Forest}}{\thesection}{.6em}{}
\titleformat{\subsection}{\large\bfseries\color{Forest}}{\thesubsection}{.55em}{}
\titleformat{\subsubsection}{\normalsize\bfseries\color{Forest}}{}{0pt}{}
\titlespacing*{\section}{0pt}{18pt plus 3pt minus 2pt}{8pt}
\titlespacing*{\subsection}{0pt}{12pt plus 2pt minus 1pt}{5pt}
\titlespacing*{\subsubsection}{0pt}{9pt}{3pt}
\setlist[itemize]{leftmargin=1.4em,itemsep=2pt,topsep=3pt,parsep=0pt}
\setlist[enumerate]{leftmargin=1.7em,itemsep=3pt,topsep=4pt,parsep=0pt}
\pagestyle{fancy}
\fancyhf{}
\fancyhead[L]{\sffamily\fontsize{9}{11}\selectfont\color{Muted}LARGE CARDINALS AT THE INCONSISTENCY FRONTIER}
\fancyhead[R]{\sffamily\fontsize{9}{11}\selectfont\color{Muted}CONTINUATION IV}
\fancyfoot[L]{\small\color{Muted}Cofinal-orbit rigidity and its consistency strength}
\fancyfoot[R]{\small\color{Muted}\thepage}
\renewcommand{\headrulewidth}{.35pt}
\renewcommand{\footrulewidth}{0pt}
\fancypagestyle{plain}{\fancyhf{}\fancyfoot[L]{\small\color{Muted}Cofinal-orbit rigidity and its consistency strength}\fancyfoot[R]{\small\color{Muted}\thepage}\renewcommand{\headrulewidth}{0pt}}
\tcbset{colback=Pale,colframe=Sage,colbacktitle=Sage,coltitle=Forest,fonttitle=\bfseries,boxrule=.5pt,arc=3pt,left=9pt,right=9pt,top=6pt,bottom=6pt,toptitle=3pt,bottomtitle=3pt,before skip=8pt,after skip=8pt,breakable}
\newtcolorbox{assessment}[1]{title={#1}}
\theoremstyle{definition}
\newtheorem{definition}{Definition}[section]
\newtheorem{theorem}[definition]{Theorem}
\newtheorem{proposition}[definition]{Proposition}
\newtheorem{lemma}[definition]{Lemma}
\newtheorem{corollary}[definition]{Corollary}
\newtheorem{remark}[definition]{Remark}
\newtheorem{question}[definition]{Question}
\newtheorem{inputthm}{Input}
\renewcommand{\theinputthm}{\Alph{inputthm}}
\numberwithin{equation}{section}
\newcommand{\ZFC}{\mathsf{ZFC}}
\newcommand{\ZF}{\mathsf{ZF}}
\newcommand{\AC}{\mathsf{AC}}
\newcommand{\GA}{\mathsf{GA}}
\newcommand{\HOD}{\mathrm{HOD}}
\newcommand{\OD}{\mathrm{OD}}
\newcommand{\CD}{\mathrm{CD}}
\newcommand{\HCD}{\mathrm{HCD}}
\newcommand{\UF}{\mathrm{UF}}
\newcommand{\Ex}{\operatorname{Ex}}
\newcommand{\UEx}{\operatorname{UEx}}
\newcommand{\Ext}{\operatorname{Ext}}
\newcommand{\SC}{\operatorname{SC}}
\newcommand{\CEx}{\operatorname{CEx}}
\newcommand{\REx}{\operatorname{REx}}
\newcommand{\Con}{\operatorname{Con}}
\newcommand{\crit}{\operatorname{crit}}
\newcommand{\cf}{\operatorname{cf}}
\newcommand{\ot}{\operatorname{ot}}
\newcommand{\ran}{\operatorname{ran}}
\newcommand{\rank}{\operatorname{rank}}
\newcommand{\lcm}{\operatorname{lcm}}
\newcommand{\Ord}{\mathrm{Ord}}
\newcommand{\Pow}{\mathcal P}
\newcommand{\id}{\mathrm{id}}
\newcommand{\restr}{\mathbin{\upharpoonright}}
\newcommand{\Izero}{\mathrm{I0}}
\newcommand{\Itwo}{\mathrm{I2}}
\newcommand{\Ithree}{\mathrm{I3}}
\newcommand{\Iwf}{\mathrm{I3}_{\mathrm{wf}}(0)}
\newcommand{\Sel}{\operatorname{Sel}}
\newcommand{\FF}{\mathcal F}
\newcommand{\PP}{\mathbb P}
\newcommand{\Clam}{\mathcal C_\lambda}
\newcommand{\Rig}{\operatorname{Rig}}
\newcommand{\Spec}{\operatorname{Spec}}
\newcommand{\ch}{\operatorname{ind}}
\newcommand{\CF}{\mathsf{CF}}
\newcommand{\src}[1]{\unskip\ {\normalfont\small\sffamily\color{Olive}[#1]}\ }
\newcommand{\R}[1]{\textsf{R#1}}
\newcolumntype{Y}{>{\raggedright\arraybackslash}X}
\newcolumntype{P}[1]{>{\raggedright\arraybackslash}p{#1}}
\newcolumntype{C}{>{\centering\arraybackslash}p{.034\textwidth}}
\newcommand{\yes}{$\bullet$}
\newcommand{\prt}{$\circ$}
\newcommand{\arxiv}[1]{\href{https://arxiv.org/abs/#1}{arXiv:#1}}
\newcommand{\doi}[1]{\href{https://doi.org/#1}{doi:\nolinkurl{#1}}}


\newcommand{\Meas}{\operatorname{Meas}}
\newcommand{\RS}{\mathsf{RS}}
\newcommand{\RC}{\mathsf{RC}}
\newcommand{\MR}{\mathsf{MR}}
\newcommand{\HR}{\mathsf{HR}}
\newcommand{\BB}{\mathbb B}
\newcommand{\AlgA}{\mathbb A}
\newcommand{\tc}{\operatorname{tc}}
\newcommand{\Fin}{\operatorname{Fin}}
\newcommand{\splus}{\mathbin{\smallfrown}}
\newcommand{\one}{\mathbf 1}
\newcommand{\symdiff}{\mathbin{\triangle}}
\newcommand{\BoolVal}[1]{\lVert #1\rVert}
\newcommand{\checklambda}{\check\lambda}
\begin{document}
\thispagestyle{empty}
{\sffamily\bfseries\small\color{Olive}SET THEORY\quad /\quad RESEARCH CONTINUATION IV}

\vspace{13pt}
{\fontsize{28}{31}\selectfont\bfseries\color{Forest}The measurable strength\\ of cofinal-orbit rigidity\par}
\vspace{10pt}
{\fontsize{14}{18}\selectfont An equiconsistency theorem and a canonical\\ weakly homogeneous forcing decomposition\par}
\vspace{14pt}
{\color{Muted}Prepared for Vladimir Reshetnikov\hfill 18 September 2026\par}
\vspace{3pt}
{\color{Sage}\hrule height .6pt}
\vspace{12pt}

\textbf{Abstract.}
The supplied synthesis proves that ultraexacting cardinals admit rigid classes of cofinal
$\omega$-sequences modulo finite difference. We determine the consistency strength of
this particular conclusion: the existence of one such class is equiconsistent with
one measurable cardinal. The same upper bound supplies $\lambda$ distinct rigid classes
with pairwise disjoint representatives, simultaneously witnessing the full-width
obstruction to definable complete sections. The upper bound is an ordinary Prikry
extension. Its essential additional feature is that the cone isomorphisms preserve
\emph{the entire finite-change class}, although they do not preserve its representative.
The lower bound is an application of the classical Dodd--Jensen covering theorem.

We also give a forcing interpretation of each rigid class $q$. For every $a\in q$,
$\HOD_{\{a\}}$ is a proper set-forcing extension of $\HOD_{\{q\}}$ by a canonically
coded, atomless, weakly homogeneous Boolean algebra. Every representative gives a
generic filter for the same algebra and the same extension. Under the additional
hypothesis $V_\lambda\subseteq\HOD$, this extension adds no bounded subsets of
$\lambda$ but changes its cofinality from $\lambda$ to $\omega$. The arguments distinguish
an orbit used as one parameter from a representative used as a parameter throughout.

\begin{assessment}{Main conclusion}
Writing $\RC$ for the rigid-class assertion defined precisely in
Section~\ref{sec:statements}, we prove
\[
 \boxed{\quad
 \Con(\ZFC+\exists\kappa\,\Meas(\kappa))
 \quad\Longleftrightarrow\quad
 \Con(\ZFC+\RC).
 \quad}
\]
This calibrates a consequence of ultraexactingness, not ultraexactingness itself.
It shows that the rigid-class package, including its full-width consequences,
already occurs at measurable-cardinal consistency strength. The new contribution is the
calibration and forcing analysis of the dossier's exact definition of rigidity;
the underlying Prikry, Vop\v enka, and Dodd--Jensen machinery is classical.
\end{assessment}

\vspace{5pt}
{\small\color{Muted}Research status. The new deductions are proved below in English.
The deep covering theorem is an explicitly identified input, not a claimed new proof
of fine structure. No global priority claim or Lean verification is asserted.}

\newpage
\tableofcontents
\newpage

\section{Statements and relation to the supplied research}\label{sec:statements}

The reference point is the second-edition \emph{Large Cardinals Synthesis} in
\nolinkurl{Cardinals3.zip}, especially its subsection ``Rigid classes'' and the
questions following its analysis of the cofinal quotient \cite{Dossier}.
We retain its definition of rigidity. We do not assume any of its proposed
consistency constructions in proving the new upper bound.

\subsection{Parameters, cofinal sets, and rigidity}

All unqualified cardinalities and all occurrences of $\OD$ are computed in the
ambient universe $V\models\ZFC$. For a collection $C$ of allowed parameters,
$\OD_C$ permits ordinal parameters and finitely many elements of $C$.
$\HOD_C$ is the corresponding hereditary class. In particular, in
$\OD_{V_\lambda\cup\{q\}}$, the set $q$ is \emph{one additional parameter};
its members are not thereby allowed as parameters. The distinction is essential.
We abbreviate $\OD_{\{q\}}$ to $\OD_q$ when no ambiguity is possible.

For an uncountable cardinal $\lambda$ of cofinality $\omega$, put
\begin{align*}
 D_\lambda&=\{a\subseteq\lambda:\ot(a)=\omega,\ \sup a=\lambda\},\\
 \Clam&=\{a\subseteq\lambda:\sup a=\lambda,\ |a|<\lambda\},\\
 a=^*b&\quad\Longleftrightarrow\quad |a\symdiff b|<\omega,\\
 Q_\lambda&=D_\lambda/{=^*}.
\end{align*}
An element of $Q_\lambda$ is a set, despite the word ``class.'' For $a\in D_\lambda$,
\[
 [a]=\{a\symdiff d:d\in[\lambda]^{<\omega}\}.
\]
Every finite modification of $a$ is still a cofinal set of order type $\omega$.
There are exactly $\lambda$ such modifications: the displayed parametrization gives
the upper bound, and $\alpha\mapsto a\symdiff\{\alpha\}$ gives the lower bound.

\begin{definition}[The dossier's rigidity property]\label{def:rig}
For $q\in Q_\lambda$, let $\Rig(\lambda,q)$ assert that
\[
 \varnothing\ne A\subseteq\Clam,\qquad
 A\in\OD_{V_\lambda\cup\{q\}}
 \quad\Longrightarrow\quad |A|\ge\lambda.
\]
The quantified family $A$ need not consist of representatives of $q$.
It can be any definable family of short cofinal subsets of $\lambda$.
\end{definition}

We will use three different hereditary classes:
\begin{equation}\label{eq:models}
 N_q=\HOD_{V_\lambda\cup\{q\}},\qquad
 H_q=\HOD_{\{q\}},\qquad
 M_a=\HOD_{\{a\}}\quad(a\in q).
\end{equation}
The first is an inner model of $\ZF$ and contains $V_\lambda$; it is not assumed
to satisfy Choice. The latter two are inner models of $\ZFC$, since a single
fixed parameter gives a canonical definable well-order of its hereditarily
definable sets. Neither this fact nor the notation implies $q\in H_q$.
Indeed, for a rigid class $q\notin H_q$.

These uses of definability are first-order expressible by the usual rank-local
coding of ordinal definitions: use a formula, a rank $V_\theta$ in which it is
evaluated, and finitely many ordinal parameters. They do not require a truth
predicate for the entire universe.

\subsection{The theories being calibrated}

\begin{definition}\label{def:theories}
Let $\RS(\lambda)$ be the conjunction
\[
 \lambda>\omega\text{ is a cardinal},\qquad
 \cf(\lambda)=\omega,\qquad
 (\forall\alpha<\lambda)\ |V_\alpha|<\lambda.
\]
This implies that $\lambda$ is strong limit and $|V_\lambda|=\lambda$.
Define the following first-order assertions over $\ZFC$.

$\HR$: there is a $\lambda$ satisfying $\RS(\lambda)$ which is regular in
$\HOD_{V_\lambda}$.

$\RC$: there are $\lambda$ and $q\in Q_\lambda$ such that
$\RS(\lambda)$ and $\Rig(\lambda,q)$.

$\MR$: there are a $\lambda$ satisfying $\RS(\lambda)$ and a sequence
$\langle a_i:i<\lambda\rangle$ of pairwise disjoint members of $D_\lambda$
such that every $q_i=[a_i]$ is rigid and all the classes $N_{q_i}$ are equal.
\end{definition}

The rank bound in $\RS$ matches the setting of the supplied reports. It is
stronger than merely saying ``strong limit of countable cofinality.'' The
uncountability requirement is also indispensable: at $\lambda=\omega$ the
family $\Clam$ is empty and the rigidity definition would be vacuous.

\begin{theorem}[Consistency calibration]\label{thm:main}
The following theories are equiconsistent:
\[
 \ZFC+\exists\kappa\,\Meas(\kappa),\qquad
 \ZFC+\HR,\qquad \ZFC+\RC,\qquad \ZFC+\MR.
\]
More specifically, $\MR\Rightarrow\RC\Rightarrow\HR$ in $\ZFC$;
$\HR$ implies the existence of an inner model with a measurable cardinal;
and ordinary Prikry forcing over a measurable cardinal forces $\MR$.
\end{theorem}

The proof occupies Sections~\ref{sec:regular}--\ref{sec:calibration}.
Sections~\ref{sec:boolean}--\ref{sec:forcing} prove the complementary forcing
structure theorem. There is no assertion that the measurable cardinal from the
lower bound is measurable in $V$, or that it must be the ambient singular
cardinal $\lambda$.

\subsection{What is extended, and what is imported}

\begin{center}
\small
\begin{tabularx}{\textwidth}{@{}P{.27\textwidth}Y@{}}
\toprule
Ingredient & Role in this continuation\\
\midrule
Rigidity from ultraexactingness & Prior dossier result; motivation only, not an input to the new measurable upper bound.\\
Rigidity versus regularity & An exact characterization, proved in Section~\ref{sec:regular}.\\
Prikry's theorem & Classical forcing machinery; the relevant proof is included. The additional parameter bookkeeping preserves the entire class $q$.\\
Dodd--Jensen covering & Deep classical input used only for the lower bound. Its application is proved explicitly.\\
Vop\v enka's algebra & Classical construction; specialized and proved here, including the common-extension and finite-change homogeneity conclusions.\\
Many copies and saturation & The underlying deductions already occur in the dossier. Here they are realized from a single measurable cardinal, rather than an ultraexacting hypothesis.\\
\bottomrule
\end{tabularx}
\end{center}

A related current result should be distinguished from this calibration.
In a preprint posted on 25 August 2026, Benhamou and coauthors
\cite{BCGHP} obtain the optimal measurable-cardinal bound for making
$\aleph_\omega$ the first ordinal at which the ambient and HOD power sets
disagree. Their statement concerns a first-disagreement phenomenon, not
the exact rigid quotient class with all $V_\lambda$ parameters used here.
It is not an input to our proof. The present contribution is a continuation
of the supplied dossier; neither this comparison nor the literature check
establishes global priority for the formulation below.

\section{Rigidity is exactly relative HOD regularity}\label{sec:regular}

The first observation removes the quantification over families from the
\emph{test} for rigidity, without weakening the property.

\begin{lemma}[Least-order-type compression]\label{lem:compression}
Suppose $\lambda$ is an infinite cardinal and
$\varnothing\ne A\subseteq\Clam$ has $|A|<\lambda$.
There is a cofinal $u\subseteq\lambda$ with $|u|<\lambda$ which is definable
from $A$ and $\lambda$, without choosing a member of $A$.
\end{lemma}
\begin{proof}
Let $\tau$ be the least ordinal occurring as the order type of a member of $A$,
and let
\[
 A_\tau=\{a\in A:\ot(a)=\tau\},\qquad u=\bigcup A_\tau.
\]
Since the members of $A$ have cardinality below the initial ordinal $\lambda$,
their order types are also below $\lambda$. Thus $\tau<\lambda$.
The family $A_\tau$ is nonempty, so $u$ is cofinal. In the ambient universe,
\[
 |u|\le |A_\tau|\cdot|\tau|<\lambda.
\]
The last inequality uses a product of two \emph{fixed} cardinals below
$\lambda$. It does not assume that $\lambda$ is regular. All the operations
used to define $u$ are canonical from $A$.
\end{proof}

\begin{lemma}[Order type detects external smallness]\label{lem:otp}
If $B\subseteq\lambda$ and $|B|^V<\lambda$, where $\lambda$ is a cardinal
of $V$, then $\ot(B)<\lambda$. If $B$ belongs to a transitive inner model,
its order type and increasing enumeration are absolute between that model and $V$.
\end{lemma}
\begin{proof}
The increasing enumeration is an order-isomorphism with a unique ordinal.
If that ordinal were at least the initial ordinal $\lambda$, it would give
an injection of $\lambda$ into $B$, contrary to the assumed cardinality.
The well-founded recursion computing the enumeration is absolute between
transitive models with the same ordinals.
\end{proof}

\begin{theorem}[Intrinsic regularity characterization]\label{thm:rig-regular}
For every uncountable cardinal $\lambda$ of cofinality $\omega$ and every
$q\in Q_\lambda$,
\[
 \Rig(\lambda,q)
 \quad\Longleftrightarrow\quad
 N_q\models\text{``$\lambda$ is regular.''}
\]
Equivalently, there is no cofinal $u\subseteq\lambda$ of order type below
$\lambda$ in $\OD_{V_\lambda\cup\{q\}}$.
\end{theorem}
\begin{proof}
If $N_q$ contains a cofinal map $f:\tau\to\lambda$ with $\tau<\lambda$,
then its range $u$ is a short cofinal set in $N_q$. In particular,
$u\in\OD_{V_\lambda\cup\{q\}}$, and the definable singleton $\{u\}$ violates
rigidity. Notice that an ordinal domain below $\lambda$ is small in $V$,
even if the two models disagree about some other cardinalities.

Conversely, suppose that rigidity fails, witnessed by a nonempty definable
$A\subseteq\Clam$ of cardinality below $\lambda$. Apply
Lemma~\ref{lem:compression} to obtain a definable short cofinal set $u$.
A definable set of ordinals is hereditarily definable from the same allowed
parameters, because every ordinal is ordinal definable. Thus $u\in N_q$.
By Lemma~\ref{lem:otp}, its increasing enumeration has an ordinal domain
$\tau<\lambda$, and the same enumeration belongs to $N_q$. This contradicts
regularity there.

Finally, a cofinal set of order type below $\lambda$ gives a singleton
counterexample, whereas compression produces such a set from any
counterexample family. This proves the equivalent formulation as well.
\end{proof}

\begin{corollary}\label{cor:hered}
If $q$ is rigid, then $\lambda$ is regular in $H_q$ and in
$\HOD_{V_\lambda}$. Moreover $q\cap N_q=\varnothing$ and $q\notin N_q$.
\end{corollary}
\begin{proof}
Both smaller hereditary classes are contained in $N_q$, so a cofinal map
in either would be one in $N_q$. A member of $q$ in $N_q$ would itself be
a cofinal set of order type $\omega<\lambda$. Finally, transitivity would
put every member of $q$ in $N_q$ if $q$ belonged to $N_q$.
\end{proof}

This characterization is useful but not a solution of every classification
question in the dossier: it does not decide which classes are rigid from
finite combinatorial data, nor does it produce distinct hereditary models
from distinct rigid classes.

\section{The measurable lower bound}\label{sec:lower}

Here is the only deep inner-model input to the calibration.

\begin{inputthm}[Dodd--Jensen covering below a measurable]\label{input:dj}
If there is no inner model of $\ZFC$ with a measurable cardinal, the
Dodd--Jensen core model $K$ has a canonical ordinal-definable well-order
and covers every uncountable set $X$ of ordinals by a set $Y\in K$ with
$X\subseteq Y$ and $|Y|^V=|X|^V$.
\end{inputthm}

We use the classical covering theorem, not a conjectural covering principle
for $\HOD$. The original sources are Dodd--Jensen \cite{DJK,DJLU};
Mitchell explicitly records the resulting singular-cardinal dichotomy in
his introduction \cite{Mitchell}. The proof of the covering theorem itself
requires core-model fine structure and is not reproduced here.

\begin{proposition}[No inner measurable, no HOD regularization]\label{prop:lower}
If an uncountable singular cardinal $\lambda$ of $V$ is regular in
$\HOD$, then there is an inner model with a measurable cardinal.
Consequently $\HR$, $\RC$, and $\MR$ each imply such an inner model.
\end{proposition}
\begin{proof}
Assume that there is no inner model with a measurable cardinal, and use
Input~\ref{input:dj}. Choose a cofinal $c\subseteq\lambda$ with
$|c|=\cf(\lambda)<\lambda$. An uncountable singular cardinal is a limit
cardinal, so $\omega_1<\lambda$. Therefore
\[
 X=c\cup\omega_1
\]
is an uncountable set of ordinals of cardinality below $\lambda$.
Cover it by $Y\in K$ of the same ambient cardinality, and put
$B=Y\cap\lambda$. Then $B\in K$, $B$ is cofinal in $\lambda$, and
$|B|^V<\lambda$.

The order type $\tau=\ot(B)$ is below $\lambda$ by
Lemma~\ref{lem:otp}. Hence $K$ contains an increasing cofinal map
$\tau\to\lambda$. Every set of $K$ is ordinal definable in $V$, by its
position in the canonical well-order of $K$; the same is true of its
transitive closure. Thus $K\subseteq\HOD$, and the cofinal map belongs to
$\HOD$. This contradicts the stipulated regularity of $\lambda$ there.

For the final assertions, regularity in $\HOD_{V_\lambda}$ implies
regularity in its submodel $\HOD$. Corollary~\ref{cor:hered} applies to
$\RC$, and $\MR$ includes an instance of $\RC$.
\end{proof}

\begin{corollary}[Conditional inconsistency]\label{cor:no-inner}
The conjunction of $\ZFC$, ``there is no inner model with a measurable
cardinal,'' and ``there is a rigid class at an uncountable singular cardinal
of cofinality $\omega$'' is inconsistent.
\end{corollary}

This is a lower bound, not an assertion that a singular cardinal becomes
measurable in the ambient universe. Nor does it use the large-cardinal
embedding assumptions from which the supplied dossier originally obtained
rigidity.

\section{A quotient-preserving homogeneity principle}\label{sec:transfer}

Ordinary homogeneity is insufficient as a slogan here: we must specify
\emph{which name} survives the cone isomorphisms.

\begin{definition}\label{def:qh}
Let $W\models\ZFC$, let $\PP\in W$ be a set forcing, and let $\dot q$
be a $\PP$-name. Say that the forcing is \emph{cone homogeneous over
$\dot q$} if, for any $p,r\in\PP$, there are $p'\le p$, $r'\le r$ and
an isomorphism between the cones below $p'$ and $r'$ which preserves
$\dot q$ and the check names for ground-model parameters.
Preservation of $\dot q$ means equality of its transported and original
interpretations under the corresponding generic filters.
\end{definition}

\begin{lemma}[The fixed-name decision lemma]\label{lem:decision}
Under Definition~\ref{def:qh}, every sentence
$\varphi(\dot q,\check x_1,\ldots,\check x_m)$, with $x_i\in W$, is
decided by the top condition. In $V=W[G]$, every set of ordinals definable
from $q=\dot q^G$ and finitely many ground-model parameters belongs to $W$.
\end{lemma}
\begin{proof}
Suppose two conditions force opposite values of the sentence. Strengthen
them to the isomorphic cones given by the hypothesis. The forcing
isomorphism theorem transports the sentence to itself, so the first cone
would force both a sentence and its negation, a contradiction. Since the
conditions deciding any sentence are dense, the top condition decides it.

Now let $A\subseteq\beta$ in the extension be defined from $q$ and ground
parameters by $\psi$. For every $\xi<\beta$, apply the first part to the
sentence expressing membership of $\xi$ in the uniquely defined set.
The forcing relation is a definable relation of $W$, and therefore
\[
 A=\{\xi<\beta: \one\Vdash\text{``$\xi$ belongs to that set''}\}
\]
is a set in $W$. The ordinal $\beta$ may be included among the parameters.
This also handles definitions originally presented as unique-object
definitions rather than as membership formulas.
\end{proof}

\begin{theorem}[Rigidity transfer]\label{thm:transfer}
Suppose $V=W[G]$ satisfies all of the following:
$\lambda$ is regular in $W$; it remains an uncountable cardinal and has
cofinality $\omega$ in $V$; $V_\lambda^V=V_\lambda^W$;
$q=\dot q^G\in Q_\lambda^V$; and $\PP$ is cone homogeneous over $\dot q$.
Then $\Rig^V(\lambda,q)$.
\end{theorem}
\begin{proof}
Assume that $A$ is a counterexample to rigidity in $V$. By
Lemma~\ref{lem:compression}, there is a cofinal $u\subseteq\lambda$,
of ambient cardinality below $\lambda$, definable from $q$, ordinal
parameters, and finitely many elements of $V_\lambda^V$.
All the latter parameters are in $W$. Lemma~\ref{lem:decision} therefore
gives $u\in W$. But $\ot(u)<\lambda$ in $V$, and its order type is
absolute to $W$. Its increasing enumeration contradicts regularity of
$\lambda$ in $W$.
\end{proof}

\begin{remark}
No claim is made that the individual representative of $q$ is fixed by the
isomorphisms. That would be incompatible with the intended singularization.
The protected information is precisely its finite-change class.
\end{remark}

\section{The Prikry construction}\label{sec:prikry}

Let $W\models\ZFC$ contain a measurable cardinal $\kappa$, and fix a
normal nonprincipal $\kappa$-complete ultrafilter $U$ on $\kappa$ in $W$.
All forcing constructions in this section are performed in $W$.
We recall enough of the ordinary Prikry argument to account for every
hypothesis of Theorem~\ref{thm:transfer}. The classical source is Prikry
\cite{Prikry}; a modern research treatment with detailed proofs is
Benhamou \cite{Benhamou}.

\subsection{Conditions and the Prikry property}

A condition in $\PP_U$ is a pair $(s,A)$, where $s$ is a finite increasing
sequence from $\kappa$, $A\in U$, and every member of $A$ is larger than
$\max(s)$ when $s\ne\varnothing$. A stronger condition $(t,B)\le(s,A)$
has $s$ as an initial segment of $t$, adds to the stem only points of $A$,
and has $B\subseteq A$. A \emph{direct} extension has the same stem.
We identify finite increasing sequences with their ranges when taking
unions or symmetric differences.

Normality has the following two familiar consequences. We include their
short derivations to make the subsequent fusion argument precise.

\begin{lemma}[Diagonalization and finite homogeneity]\label{lem:normal}
For a family $\langle A_t:t\in[\kappa]^{<\omega}\rangle$ of members of $U$,
there is $B\in U$ such that
\[
 t\in[B]^{<\omega},\ \beta\in B,\ \max(t)<\beta
 \quad\Longrightarrow\quad \beta\in A_t.
\]
The empty-stem requirement can be imposed by intersecting with
$A_\varnothing$. Also, if $h:[\kappa]^{<\omega}\to r$ with $r$ finite,
there is $B\in U$ on which $h$ is constant separately on each
$[B]^n$, $n<\omega$.
\end{lemma}
\begin{proof}
For each $\alpha<\kappa$, intersect the $A_t$ for $t\subseteq\alpha$.
There are fewer than $\kappa$ such finite sets, so the intersection is in $U$.
Take the normal diagonal intersection of these intersections and restrict
it to limit ordinals. Given $\max(t)<\beta$ in this diagonal set, choose
$\alpha$ strictly between $\max(t)$ and $\beta$; the desired membership
follows. For empty $t$, make the indicated extra intersection. Normal
measures contain all clubs, including the set of limit ordinals.

For the second assertion, first prove homogeneity on $n$-element sets by
induction on $n$. The case $n=1$ is the ultrafilter property. For a coloring
of $(n+1)$-element sets, and each $n$-element stem $t$, choose the unique
color taken on a $U$-large set of possible last points. Diagonalize these
large sets by the first assertion. Apply the inductive hypothesis to the
resulting coloring of $n$-element stems, and intersect with the diagonal
set. This makes the original coloring constant on $(n+1)$-element sets.
Finally intersect the homogeneous sets obtained for all finite $n$.
Countable completeness preserves membership in $U$.
\end{proof}

\begin{lemma}[Prikry property]\label{lem:prikry}
For any forcing sentence $\varphi$ and any $(s,A)\in\PP_U$, a direct
extension of $(s,A)$ decides $\varphi$.
\end{lemma}
\begin{proof}
For each finite extension $t$ of the stem $s$, assign a color $h(t)$:
use $1$ if some condition with stem $t$ forces $\varphi$, use $-1$ if
some condition with that stem forces $\neg\varphi$, and use $0$ if no
condition with that stem decides the sentence. The first two alternatives
cannot both occur, because two conditions with the same stem are
compatible by intersection of their measure-one sets. In a nonzero case,
choose a witnessing tail set $A_t\in U$; in a zero-color case put $A_t=\kappa$.

Apply Lemma~\ref{lem:normal} to the coloring of finite continuations $u$
given by $h(s\splus u)$. Then diagonalize the chosen tail sets, retaining
a measure-one subset $B\subseteq A$ such that the tail of $B$ after $u$
is included in $A_{s\splus u}$. The requirements for the fixed stem $s$
and for the empty continuation are met by an additional intersection.
On $B$, the color of $s\splus u$ depends only on the length of $u$.

Some extension of $(s,B)$ decides $\varphi$, by density. Let $n$ be its
number of additional stem points and let $\varepsilon\in\{1,-1\}$ be
its color. For every $v\in[B]^n$, the condition
\[
 (s\splus v,\ B\setminus(\max(v)+1))
\]
forces the same value $\varepsilon$, using the diagonal refinement.
For $n=0$, read this as the condition with stem $s$ itself. Every condition
below $(s,B)$ can be extended to have at least $n$ additional stem points;
it then extends one of these deciding conditions. Thus a dense set below
$(s,B)$ forces the same value, and $(s,B)$ decides $\varphi$.
\end{proof}

\subsection{The extension preserves the entire lower rank}

\begin{proposition}\label{prop:prikry-facts}
If $G\subseteq\PP_U$ is generic over $W$ and $V=W[G]$, then all cardinals
are preserved, $\cf^V(\kappa)=\omega$, no bounded subsets of $\kappa$ are
added, and
\[
 V_\kappa^V=V_\kappa^W,\qquad
 (\forall\alpha<\kappa)\ |V_\alpha|^V<\kappa.
\]
\end{proposition}
\begin{proof}
The union $c$ of the stems in $G$ is an increasing sequence of length
$\omega$. Conditions with arbitrarily long stems are dense, and for each
$\beta<\kappa$ the conditions placing a stem point above $\beta$ are dense.
Hence $c$ is cofinal in $\kappa$.

Let $\mu<\kappa$, and let $\dot x$ be a name for a subset of $\mu$.
Starting from any condition, decide the statements $\xi\in\dot x$,
$\xi<\mu$, successively by direct extensions using
Lemma~\ref{lem:prikry}. At limits, intersect the tail sets.
$\kappa$-completeness supplies a single direct extension deciding the entire
subset. Therefore every subset of $\mu$ in $V$ already belongs to $W$.
In particular, every cardinal below $\kappa$ is preserved, since a
bijection witnessing its collapse could be coded by a bounded subset of
$\kappa$.

The cardinal $\kappa$ is inaccessible in $W$. Its smaller cardinals remain
cardinals and are unbounded in $\kappa$, so $\kappa$ cannot be collapsed:
a purported injection into a smaller ordinal would already collapse some
preserved cardinal below $\kappa$. Finally, there are only $\kappa$ finite
stems, and conditions with the same stem are compatible. Thus $\PP_U$
has the $\kappa^+$-chain condition and preserves all cardinals above
$\kappa$ as well.

For the rank assertion, argue by induction on $\alpha<\kappa$.
In $W$, every $V_\alpha$ has cardinality below $\kappa$. Once $V_\alpha$
has been shown unchanged, a ground-model bijection with a smaller ordinal
codes every subset of $V_\alpha$ by a bounded subset of $\kappa$.
Consequently its power set has not changed either. Limit stages follow
by taking unions. The ground-model bijections witnessing
$|V_\alpha|<\kappa$ are still available in the extension.
\end{proof}

\subsection{The tail class, unlike the sequence, is preserved}

Let $\dot c$ be the canonical generic cofinal set and define the name
\[
 \dot q=\{\dot c\symdiff\check d:d\in[\kappa]^{<\omega}\}.
\]
In every generic extension, $\dot q$ is exactly $[\dot c]$.

\begin{lemma}[Prikry cones preserve the finite-change class]\label{lem:prikry-q}
$\PP_U$ is cone homogeneous over $\dot q$.
\end{lemma}
\begin{proof}
Take $p=(s,A)$ and $r=(t,B)$. Choose
\[
 C\in U,\qquad C\subseteq A\cap B,
 \qquad \min C>\max(s\cup t),
\]
with the evident interpretation for empty stems. Put $p'=(s,C)$ and
$r'=(t,C)$. The map
\[
 (s\splus u,D)\longmapsto(t\splus u,D)
\]
is an order isomorphism between the two cones. The same finite continuation
$u$ and tail set $D$ are permitted on either side. The stems $s$ and $t$
need not have equal lengths.

Corresponding generics have cofinal sets $s\cup h$ and $t\cup h$ for a
common tail $h$. These differ by the finite set $s\symdiff t$, so their
finite-change classes are equal. Moreover each cofinal set is definable
from the other and the two finite ground-model stems; the resulting
forcing extensions are the same. Thus the transported name for $q$ has
exactly the required interpretation, and check names for ground-model
parameters are unchanged.
\end{proof}

\begin{theorem}[A rigid class from one measurable]\label{thm:prikry-rig}
In every ordinary Prikry extension $V=W[G]$ as above,
\[
 \RS(\kappa)\quad\text{and}\quad\Rig^V(\kappa,[c])
\]
hold. In particular,
$\HOD^V_{V_\kappa\cup\{[c]\}}$ regards $\kappa$ as regular.
\end{theorem}
\begin{proof}
Proposition~\ref{prop:prikry-facts} gives $\RS(\kappa)$ and all rank and
cardinality hypotheses of Theorem~\ref{thm:transfer}.
Lemma~\ref{lem:prikry-q} gives the remaining homogeneity hypothesis.
Apply the transfer theorem, and then Theorem~\ref{thm:rig-regular}.
\end{proof}

\begin{remark}[Uniformity over low-rank parameters]
This conclusion is not restricted to parameters that were ordinal definable
in $W$. Every member of the entire ambient $V_\kappa$ is a ground-model
parameter, and the cone isomorphisms preserve all their check names.
This is why the construction proves the dossier's full
$\OD_{V_\kappa\cup\{q\}}$ version of rigidity.
\end{remark}

\section{Multiplicity, saturation, and the completed calibration}\label{sec:calibration}

\subsection{One class supplies many copies without additional strength}

\begin{lemma}[Disjoint mutually definable copies]\label{lem:copies}
For every $q\in Q_\lambda$ there are $\lambda$ distinct classes
$q_i\in Q_\lambda$, $i<\lambda$, mutually ordinal definable with $q$,
having pairwise disjoint representatives. Thus if $q$ is rigid, every
$q_i$ is rigid and $N_{q_i}=N_q$.
\end{lemma}
\begin{proof}
Use the canonical pairing $\pi:\lambda\times\lambda\to\lambda$ obtained
by ordering pairs first by their maximum coordinate and then
lexicographically. This well-order has order type $\lambda$: every pair
has fewer than $\lambda$ predecessors, and there are $\lambda$ pairs.
For each $i<\lambda$, the map $f_i(\alpha)=\pi(i,\alpha)$ is strictly
increasing and cofinal in $\lambda$. Its range is disjoint from the ranges
of the other $f_j$.

Choose any $a\in q$ and put $a_i=f_i``a$, $q_i=[a_i]$.
These representatives are pairwise disjoint cofinal $\omega$-sets.
The definition of $q_i$ does not depend on the chosen $a$: replacing $a$
by a finite modification only finitely modifies $f_i``a$. Thus $q_i$ is
definable from $q$ and $i$.

Conversely, for any $b\in q_i$, pull $b\cap\ran(f_i)$ back along $f_i$.
The resulting set differs finitely from $a$, regardless of $b$.
Its equivalence class is therefore $q$, giving a definition of $q$ from
$q_i$ and $i$. Substitution of these definitions gives equality of the
corresponding $\OD$ classes, with or without the allowed $V_\lambda$
parameters. Their hereditary subclasses are consequently equal.
\end{proof}

The copying principle is already present in the supplied synthesis. Its
role here is to show that the multiplicity appearing there costs no more
consistency strength than the first rigid class. These copies do \emph{not}
answer the question whether there are distinct models $N_q$.

\subsection{The definable-section obstruction already occurs in this model}

\begin{definition}
A \emph{complete section} is a set $T\subseteq D_\lambda$ meeting every
$q\in Q_\lambda$. It is \emph{thin} if every such intersection has
cardinality below $\lambda$. A partial section need not meet every class;
it is thin when all its nonempty fibres are smaller than $\lambda$.
\end{definition}

\begin{proposition}[Simultaneous full-width fibres]\label{prop:width}
If $q$ is rigid and $q_i$ are its copies from Lemma~\ref{lem:copies}, then
for every $T\subseteq D_\lambda$ in $\OD_{V_\lambda}$ and every $i<\lambda$,
\[
 |T\cap q_i|\in\{0,\lambda\}.
\]
The same $\lambda$ classes serve all these sets $T$ simultaneously.
Every definable complete section therefore has $\lambda$ full-sized fibres,
and no thin complete section is in $\OD_{V_\lambda}$.
\end{proposition}
\begin{proof}
$T\cap q_i$ is definable from $q_i$ and allowed low-rank parameters.
If nonempty, it is a family of short cofinal subsets of $\lambda$, so
rigidity gives cardinality at least $\lambda$. It is a subset of $q_i$,
whose cardinality is $\lambda$. No choice of $q_i$ in this argument depends
on $T$.
\end{proof}

\begin{proposition}[Small families do not evade the obstruction]\label{prop:smallfamily}
If $q$ is rigid and $\mathcal T\in\OD_{V_\lambda}$ has
$|\mathcal T|<\lambda$, with each $T\in\mathcal T$ a thin partial section,
then $(\bigcup\mathcal T)\cap q=\varnothing$.
Consequently no nonempty such family consists of complete sections,
and no such family of partial sections has complete union.
\end{proposition}
\begin{proof}
For each $T\in\mathcal T$ meeting $q$, put
$u_T=\bigcup(T\cap q)$. All members of $q$ are countable, so
$|u_T|\le |T\cap q|\cdot\aleph_0<\lambda$, and $u_T$ is cofinal.
The family of these $u_T$ is definable from $q$ and $\mathcal T$ and has
size at most $|\mathcal T|<\lambda$. Rigidity says it must be empty.
Notice that we do not take the union of the possibly unbounded sizes
$|T\cap q|$ over all $T$; the family of compressed cofinal sets is the
counterexample directly.
\end{proof}

These consequences were obtained from ultraexactingness in the dossier.
The new point is that Theorem~\ref{thm:prikry-rig} and
Lemma~\ref{lem:copies} realize all of them in an ordinary measurable-cardinal
Prikry extension.

\subsection{Proof of the equiconsistency theorem}

\begin{proof}[Proof of Theorem~\ref{thm:main}]
$\MR\Rightarrow\RC$ follows by taking $i=0$. The implication
$\RC\Rightarrow\HR$ is Corollary~\ref{cor:hered}.
Proposition~\ref{prop:lower} gives the measurable inner-model lower bound
from $\HR$, and hence the corresponding consistency implication.

For the other direction, begin with a model with a measurable cardinal,
normalize a witnessing measure, and use ordinary Prikry forcing.
Theorem~\ref{thm:prikry-rig} supplies a $\kappa$ satisfying $\RS$ and a
rigid class at $\kappa$. Lemma~\ref{lem:copies} upgrades this to $\MR$.
Thus the consistency of one measurable cardinal gives the consistency of
the strongest of the three calibrated assertions. Together the two
directions give all the stated equiconsistencies.

As usual, this is the formal relative-consistency argument using the
forcing theorem and the core-model theorem within models of set theory.
It does not add the assumption that there exists a transitive \emph{set}
model of either theory. Presenting the construction with transitive
$W\subseteq W[G]$ is a convenient semantic notation, not a stronger
hypothesis in the consistency statement.
\end{proof}

\begin{assessment}{Consequence for the inconsistency search}
The existence of rigid classes, the $0$-or-$\lambda$ law for definable
fibres, and the exclusion of small definable families of thin sections
are compatible with $\ZFC$ at exactly measurable-cardinal consistency
strength. Any contradiction obtained from this package alone would also refute the
consistency of a measurable cardinal. To go beyond that calibration, an
argument must use additional structure: for example the embedding dynamics, simultaneous preservation
of predicates, or interaction with large cardinals below $\lambda$.
\end{assessment}

\section{The canonical Boolean algebra of a class}\label{sec:boolean}

We now extract a second result from rigidity. This construction is a
specialization of Vop\v enka's theorem. A modern statement with the
single-parameter generality needed here is Goldberg's Theorem~4.8
\cite{Goldberg}. We give the construction and the extension-equality
argument, so no step is hidden in the word ``generic.''

Throughout this section, $q=[a]\in Q_\lambda$; rigidity will be invoked
only where explicitly stated. Define, externally in $V$,
\[
 \AlgA_q=\{A\subseteq q:A\in\OD_q\}.
\]
This is a set of subsets of $q$, with the ordinary Boolean operations.
It need not belong to $H_q$. Choose its canonical $\OD_q$ enumeration
\[
 f:\gamma\longrightarrow\AlgA_q
\]
by the least codes for ordinal definitions from $q$.
Transfer the Boolean operations to the ordinal $\gamma$ using $f$;
write $\BB_q$ for the resulting coded algebra.

\begin{lemma}[Internal completeness]\label{lem:complete}
$\BB_q$ is a complete Boolean algebra in $H_q$.
The decoding map $f$ is used externally and is not asserted to belong to $H_q$.
\end{lemma}
\begin{proof}
The code for the algebra consists of relations and functions on ordinals
which are ordinal definable from $q$. Their transitive closures consist
of ordinal data and such definable graphs. They therefore belong to $H_q$.
The Boolean-algebra identities are inherited from subsets of $q$.

For a family $I\subseteq\gamma$ belonging to $H_q$, the set $I$ is itself
$\OD_q$. Hence
\[
 \bigcup\{f(\xi):\xi\in I\}
\]
is an $\OD_q$ subset of $q$, and has a code in $\BB_q$. That code is the
supremum of $I$. This verifies completeness for all families which
\emph{belong to $H_q$}, exactly as required. We are not claiming external
completeness for arbitrary families in $V$.
\end{proof}

\begin{definition}
For each $b\in q$, let
\[
 G_b=\{\xi<\gamma:b\in f(\xi)\}.
\]
We view $G_b$ as an ultrafilter on $\BB_q$.
\end{definition}

\begin{lemma}[Every representative gives a generic filter]\label{lem:pointgeneric}
$G_b$ is $H_q$-generic for $\BB_q$ for every $b\in q$.
Furthermore $b\in H_q[G_b]$ and $G_b\in M_b$.
\end{lemma}
\begin{proof}
Membership at a point respects intersections and complements, so $G_b$
is an ultrafilter. Let $\mathcal A\in H_q$ be a maximal antichain of
positive elements of $\BB_q$. Its decoded sets are pairwise disjoint.
Their union is $q$: otherwise the code for the nonempty complement would
be a nonzero element incompatible with every member of $\mathcal A$,
contradicting maximality. Thus exactly one decoded member contains $b$,
and $G_b$ meets $\mathcal A$. This proves genericity.

For $\alpha<\lambda$, let $e_\alpha$ be the code for
$\{x\in q:\alpha\in x\}$. The sequence
$\langle e_\alpha:\alpha<\lambda\rangle$ is ordinal coded and $\OD_q$,
so belongs to $H_q$. Then
\[
 b=\{\alpha<\lambda:e_\alpha\in G_b\},
\]
which reconstructs $b$ in the generic extension.

Finally, $q$ is ordinal definable from $b$, as its finite-change class.
Thus $G_b$ is ordinal definable from $b$, and is a set of ordinals.
It follows that $G_b\in M_b$.
\end{proof}

\subsection{Why the extension is exactly the representative HOD}

\begin{theorem}[Exact generic-extension identity]\label{thm:vopenka}
For every $b\in q$,
\[
 H_q[G_b]=\HOD_{\{q,b\}}=M_b.
\]
For any $a,b\in q$, $M_a=M_b$. We may consequently denote this common
extension by $M_q$.
\end{theorem}
\begin{proof}
Because $q$ is ordinal definable from $b$, adjoining $q$ as an allowed
parameter in addition to $b$ does not change the hereditary class:
$\HOD_{\{q,b\}}=M_b$. Also $H_q\subseteq M_b$ and $G_b\in M_b$.
The transitive model $M_b$ evaluates all $H_q$-names, so
$H_q[G_b]\subseteq M_b$.

For completeness, we prove the reverse inclusion by explicitly coding
hereditarily definable sets. There is a fixed definable enumeration
scheme $F_z(\xi)$, $\xi\in\Ord$, of the sets ordinal definable from
$q,z$. To obtain it, well-order the rank-local definition codes and use
$\varnothing$ for failed uniqueness codes. Every $\OD_{q,z}$ set occurs.
The evaluation of each code is a first-order set-theoretic operation;
no global truth predicate is involved.

Fix $x\in M_b$. Choose an ordinal $\eta$ large enough that every member
of $\tc(\{x\})$ occurs among $F_b(\xi)$ for $\xi<\eta$. Such a bound
exists because the transitive closure is a set and every one of its
elements is ordinal definable from $b$. Increase $\eta$ if necessary to
include a code $\xi_x$ for $x$ itself.

For each $z\in q$ define
\[
 D_z=\{\xi<\eta:\tc(\{F_z(\xi)\})
                  \subseteq\ran(F_z\restr\eta)\}.
\]
The image $F_z``D_z$ is transitive. Indeed, if a value has its full
transitive closure represented below $\eta$, so does each of its
members. Keep only the least index for each repeated value, obtaining
$D_z^*\subseteq D_z$. Put a relation on this set by
\[
 \xi\ E_z\ \nu\quad\Longleftrightarrow\quad F_z(\xi)\in F_z(\nu).
\]
This is a well-founded extensional relation isomorphic to the transitive
set $F_z``D_z$.

For each $\xi,\nu<\eta$, the subsets of $q$ expressing
$\xi\in D_z^*$ and $\xi\ E_z\ \nu$ are $\OD_q$.
Their Boolean codes, indexed by $\eta$ and $\eta\times\eta$, form sets
of ordinal data in $H_q$. Evaluating them with $G_b$ therefore constructs
$D_b^*$ and $E_b$ inside $H_q[G_b]$. This relation is well-founded in
$V$, hence also in the inner model $H_q[G_b]$, and its Mostowski collapse
is absolute. The collapse at the least index representing $x$ is exactly
$x$. That index is an ordinal, so it is available in $H_q[G_b]$.
Thus $x\in H_q[G_b]$.

Finally, if $a=^*b$, the finite set $a\symdiff b$ is ordinal definable
using its finitely many ordinal elements. Each of $a$ and $b$ is
therefore ordinal definable from the other. Substitution of definitions,
including those for all elements of the transitive closure, gives
$M_a=M_b$.
\end{proof}

The identity is stronger than merely observing that the representatives
are generic over some inner model. It identifies a \emph{single} forcing
algebra, a \emph{single} ground, and a \emph{single} extension shared by
all representatives of the class.

\section{Homogeneity, singularization, and rank agreement}\label{sec:forcing}

\subsection{Finite changes give an internal automorphism action}

Let $\Gamma=([\lambda]^{<\omega},\symdiff)$. This is an ordinal-definable
group and belongs to $H_q$. For $d\in\Gamma$ define
\[
 \rho_d(A)=\{x\symdiff d:x\in A\}\qquad(A\in\AlgA_q).
\]
Transporting $\rho_d$ through $f$ gives an automorphism of $\BB_q$.
The entire coded action is in $H_q$.

\begin{lemma}[Weak homogeneity]\label{lem:weak}
The action of $\Gamma$ makes $\BB_q$ weakly homogeneous in $H_q$:
for any nonzero $p,r\in\BB_q$, some $d\in\Gamma$ has
$\rho_d(p)\wedge r\ne0$. Moreover
\[
 \rho_d[G_b]=G_{b\symdiff d}.
\]
\end{lemma}
\begin{proof}
Choose $x\in f(p)$ and $y\in f(r)$. Both lie in $q$, so
$d=x\symdiff y$ is finite. The point $y$ lies in
$\rho_d(f(p))\cap f(r)$, proving compatibility. The witness $d$ is an
element of the internal group $\Gamma$, so this establishes precisely
the asserted statement of $H_q$. The equality of filters follows by
applying the equivalence
$b\in A\Longleftrightarrow b\symdiff d\in\rho_d(A)$
to each Boolean code.
\end{proof}

\begin{lemma}[Atomlessness under rigidity]\label{lem:atomless}
If $q$ is rigid, $\BB_q$ is atomless in $H_q$.
\end{lemma}
\begin{proof}
A nonzero code decodes to a nonempty $\OD_q$ subset $A$ of $q$.
Since $A\subseteq\Clam$, rigidity gives $|A|\ge\lambda$.
Choose distinct $x,y\in A$ and an ordinal $\alpha\in x\symdiff y$.
Intersecting $A$ with $\{z\in q:\alpha\in z\}$ and with its complement
splits its code into two nonzero elements. All splitting codes are in
$H_q$. Thus no nonzero element is an atom.
\end{proof}

\begin{theorem}[Canonical singularization attached to rigidity]\label{thm:forcing-main}
Suppose $\RS(\lambda)$ and $\Rig(\lambda,q)$. There is a complete Boolean
algebra $\BB_q\in H_q$ such that:
\begin{enumerate}
\item $\BB_q$ is atomless and weakly homogeneous, with the finite-change
group acting as above.
\item Every $b\in q$ yields an $H_q$-generic $G_b$, and
$H_q[G_b]=M_q$ independently of $b$.
\item $H_q$ regards $\lambda$ as inaccessible; $M_q$ regards it as a
strong limit cardinal of cofinality $\omega$. In particular $H_q\ne M_q$.
\item $H_q$ satisfies that $\BB_q$ forces $\lambda$ to remain a cardinal
and to have cofinality $\omega$. The algebra does not have the
$\lambda$-chain condition in $H_q$.
\end{enumerate}
If additionally $V_\lambda\subseteq\HOD^V$, then $H_q=N_q$,
\[
 V_\lambda^{H_q}=V_\lambda^{M_q}=V_\lambda^V,
\]
and $\BB_q$ forces that no bounded subsets of $\lambda$ are added.
\end{theorem}
\begin{proof}
The first two clauses are Lemmas~\ref{lem:complete},
\ref{lem:pointgeneric}, \ref{lem:weak}, \ref{lem:atomless} and
Theorem~\ref{thm:vopenka}. Corollary~\ref{cor:hered} gives regularity in
$H_q$. The ordinal $\lambda$ is a cardinal of each inner model, since
any inner-model collapse would be an ambient collapse.

Strong-limitness is downward absolute here. If an inner model of Choice
had an injection of $\lambda$ into its power set of some $\mu<\lambda$,
that would be an ambient injection into $\Pow^V(\mu)$, contradicting the
ambient strong-limit property. Thus $\lambda$ is strong limit in $H_q$
and in $M_q$. It is inaccessible in the former. In the latter any
$b\in q$ supplies its cofinal $\omega$-enumeration. That enumeration
cannot belong to $H_q$, so the extension is proper.

In a weakly homogeneous Boolean algebra, the Boolean value of a sentence
with check parameters from the ground is either zero or one. Indeed,
if nonzero conditions forced opposite answers, an automorphism could
make their images compatible, while fixing those check parameters.
The assertion that $\lambda$ is a cardinal of cofinality $\omega$ holds
in the actual extension $M_q$. The truth lemma and weak homogeneity
therefore imply that it is forced by the top condition in $H_q$.

If $\BB_q$ had the $\lambda$-chain condition, take a name for a cofinal
map $\omega\to\lambda$. For each $n$, a maximal antichain deciding its
$n$th value has size below $\lambda$. Regularity of $\lambda$ in $H_q$
then bounds the union of all possible values for all $n<\omega$ below
$\lambda$, contradicting cofinality. Hence the chain condition fails.

Now assume $V_\lambda\subseteq\HOD^V$. Each low-rank parameter can be
replaced in a definition by an ordinal definition. Consequently
$\OD_{V_\lambda\cup\{q\}}=\OD_q$, so $N_q=H_q$.
Both $H_q$ and $M_q$ contain the full ambient $V_\lambda$ and are
contained in $V$, proving the rank equality. In particular the actual
extension has no new bounded subsets of $\lambda$.

To obtain the forcing assertion rather than just a statement about this
one generic extension, let $X=(V_\lambda)^{H_q}\in H_q$. The sentence
that every subset of every $\alpha<\lambda$ belongs to $\check X$
uses only ground check parameters. It holds in $M_q$, so the preceding
weak-homogeneity argument makes its Boolean value one.
\end{proof}

\begin{assessment}{Three distinctions that must not be suppressed}
First, $H_q$ is a ground of $M_q$, not necessarily a ground of the whole
ambient universe $V$. Second, preserving $\lambda$ as a cardinal does
not assert preservation of every cardinal above it. Third, adding no
bounded subsets of $\lambda$ does \emph{not} mean
$\omega$-distributivity: the extension adds a new countable cofinal
sequence of ordinals below $\lambda$.
\end{assessment}

There is also no forcing assertion of the form
$\dot b\in\check q$ here. For a rigid class, $q\notin H_q$,
so that is not a legitimate use of a ground-model check name.
The class is used externally to construct the ordinal-coded algebra;
its actual representatives then induce the generic filters.

\subsection{The representatives cannot be mutually generic}

\begin{proposition}\label{prop:not-mutual}
If $q$ is rigid and $a,b\in q$, the filters $G_a$ and $G_b$ are not
mutually generic over $H_q$ for $\BB_q\times\BB_q$.
There are exactly $\lambda$ distinct point filters $G_b$, $b\in q$,
and all of them generate $M_q$.
\end{proposition}
\begin{proof}
Different points give different filters by the membership codes
$e_\alpha$, so there are $\lambda$ of them. Their extensions coincide
by Theorem~\ref{thm:vopenka}. They are proper because
Corollary~\ref{cor:hered} makes $\lambda$ regular in $H_q$, whereas
each $b\in q$ gives a cofinal $\omega$-sequence in $M_q$.

Recall the intersection fact for mutually generic extensions:
if $G,H$ are mutually generic over a transitive ground $N$, then
$N[G]\cap N[H]=N$. Here is the relevant proof.
For a common subset $x$ of an ordinal, take names $\sigma,\tau$ from
the two coordinates and a product condition $(p,r)$ forcing their equality.
If two extensions of $p$ forced opposite answers to
$\alpha\in\sigma$, first extend $r$ to decide $\alpha\in\tau$.
One of the two resulting product conditions would contradict the forced
equality. Thus $p$ decides every membership question, making $x$ a
ground-model set. For an arbitrary common set outside $N$, choose one
of least rank. Its members are all in $N$; a ground well-order of the
appropriate ground rank converts it to a common set of ordinals outside
$N$, contradicting the case just proved.

Mutual genericity of $G_a,G_b$ would consequently make their common
extension equal to $H_q$, contrary to properness.
\end{proof}

This does not say that the algebra has no product-generic filters in a
further universe. It says that the representatives of this one ambient
finite-change class cannot form such a pair.

\section{What the result does and does not settle}\label{sec:scope}

The consistency theorem identifies a precise boundary inside the supplied
research program. Rigidity has exactly measurable consistency strength,
and its definable-section consequences already occur in the same construction. The ultraexacting embedding hypothesis is not
needed to realize them, and its additional force must be sought in
properties not used by the Prikry construction.

Theorem~\ref{thm:forcing-main} adds a second boundary. Under
$V_\lambda\subseteq\HOD$, the ``regular inside, singular outside'' picture
from the synthesis is realized by an actual weakly homogeneous set forcing
\emph{between the two indicated hereditary models}. This converts a
cofinality discrepancy into a specific forcing decomposition. It neither
identifies that forcing with ordinary Prikry forcing nor derives a normal
measure in $H_q$. The measurable lower bound is an inner-model conclusion,
not a classification of $\BB_q$.

\subsection{Remaining questions with sharper targets}

\begin{question}
Must there be distinct rigid $q,r$ with $N_q\ne N_r$ under an
ultraexacting hypothesis? The measurable construction and the copying
lemma deliberately give many classes with the \emph{same} hereditary
model, so multiplicity of classes alone cannot answer this.
\end{question}

\begin{question}
What extra property of $\BB_q$, arising from an ultraexacting embedding
rather than from rigidity alone, distinguishes it from the measurable
Prikry examples? Potential properties must involve the embedding action
or a simultaneous family of predicates; atomlessness, weak homogeneity,
and singularization by themselves are insufficient.
\end{question}

\begin{question}
Does ordinary exactingness imply the existence of a rigid class, or the
nonexistence of a definable thin complete section, in full generality?
The consistency calibration does not prove this implication. The
quotient-preserving cone maps used in the upper bound belong to a
particular forcing construction, not to every exacting witness.
\end{question}

We make no preservation claim for a strongly compact cardinal below the
Prikry point, no new consistency claim for cover-exactingness above strong
compactness, and no determination of the successor-collapse alternative
from the synthesis. Those are different targets and require additional
arguments. The published consistency relationship between ultraexactingness
and $\Izero$ likewise remains a separate result \cite{ABGL}.

\section{Proof audit and dependency ledger}\label{sec:audit}

\begin{center}
\small
\begin{tabularx}{\textwidth}{@{}P{.28\textwidth}Y@{}}
\toprule
Conclusion & Logical dependencies\\
\midrule
Rigidity $\Leftrightarrow$ regularity of $N_q$ & Definition of rigidity; least-order-type compression; absoluteness of ordinal order type. No large-cardinal input.\\
Measurable lower bound & Classical Dodd--Jensen covering and its canonical core-model well-order; the application in Proposition~\ref{prop:lower}.\\
Measurable upper bound & A normal measure; the Prikry property and rank preservation proved here; cone isomorphisms preserving the orbit name.\\
$\lambda$ simultaneous rigid copies & Definable ordinal pairing and substitutions of ordinal definitions. This copying mechanism is also in the dossier.\\
Exact forcing decomposition & Ordinal-coded Vop\v enka algebra; maximal-antichain genericity; the explicit collapse-of-definition-codes proof.\\
Weak homogeneity and atomlessness & Finite symmetric differences; rigidity; ordinal membership tests.\\
No bounded subsets in the canonical extension & The \emph{additional} assumption $V_\lambda\subseteq\HOD^V$, followed by weak homogeneity. It is not inferred from rigidity alone.\\
\bottomrule
\end{tabularx}
\end{center}

\subsection{Checks against common proof failures}

The cardinal bounds are ambient bounds until a model superscript is given.
When transferring smallness to an inner model, the proof uses the
\emph{absolute order type} of a subset of $\lambda$, not an unjustified
assumption of cardinal correctness. The compression proof does not union
an unbounded collection of sizes below a singular cardinal. The Prikry
homogeneity argument fixes the class name and all ground check parameters,
not the generic sequence. The Boolean algebra is complete in $H_q$, not
necessarily in $V$. Its ordinal code belongs to $H_q$, while its decoding
map and the class $q$ need not. Finally, the assertion about all generics
adding no bounded subsets is deduced from weak homogeneity and ground
check parameters, not from the behavior of one extension alone.

\subsection{Verification status}

All newly used deductions are written as mathematical proofs in this
report. The supplied Lean files were reference material for the definitions
and the scope of the earlier claims; the present equiconsistency argument
and the canonical forcing analysis have not been formalized in Lean.
The Dodd--Jensen theorem is an imported classical theorem. The report does
not claim an independent formal check or a new proof of that theorem.
Bibliographical verification establishes the cited classical inputs, not
global novelty of the particular formulation proved here.

\begin{thebibliography}{99}

\bibitem{Dossier}
\emph{Large cardinals at the inconsistency frontier: a synthesis}.
Second-edition research report supplied by the user in
\nolinkurl{Cardinals3.zip},
\nolinkurl{docs/Large_Cardinals_Synthesis.tex}, 18 September 2026.
In particular the subsection ``Rigid classes,'' its corollaries and
sharpness remarks. The companion
\nolinkurl{docs/Large_Cardinals_Unified_Report.tex} supplies the broader
research program.

\bibitem{DJK}
A.~J. Dodd and R.~B. Jensen.
\emph{The covering lemma for $K$}.
Annals of Mathematical Logic \textbf{22} (1982), no.~1, 1--30.
Classical core-model input to the lower bound.

\bibitem{DJLU}
A.~J. Dodd and R.~B. Jensen.
\emph{The covering lemma for $L[U]$}.
Annals of Mathematical Logic \textbf{22} (1982), no.~2, 127--135.

\bibitem{Mitchell}
William~J. Mitchell.
\emph{Definable singularity}.
Transactions of the American Mathematical Society \textbf{327} (1991),
no.~1, 407--426.
\doi{10.1090/S0002-9947-1991-1036006-0}.
The introduction, especially p.~408, records the Dodd--Jensen
singular-cardinal alternatives used to check the lower-bound input.
\url{https://www.ams.org/journals/tran/1991-327-01/S0002-9947-1991-1036006-0/S0002-9947-1991-1036006-0.pdf}.

\bibitem{Prikry}
Karel Prikry.
\emph{Changing measurable into accessible cardinals}.
Dissertationes Mathematicae \textbf{68} (1970), 55 pp.
The original normal-measure singularization construction.

\bibitem{Benhamou}
Tom Benhamou.
\emph{Prikry forcing and tree Prikry forcing of various filters}.
\arxiv{1801.04424}, version~2, 2018.
Section~3 gives the standard Prikry facts; Corollary~3.24 records the
finite-change criterion for equality of ordinary Prikry extensions.
\url{https://arxiv.org/html/1801.04424v2}.

\bibitem{KRS}
Peter Koepke, Karen R\"asch, and Philipp Schlicht.
\emph{A minimal Prikry-type forcing for singularizing a measurable cardinal}.
Journal of Symbolic Logic \textbf{78} (2013), no.~1, 85--100.
The introduction also discusses the Dodd--Jensen lower bound for
singularization by forcing. Its minimal forcing is \emph{not} the
ordinary Prikry forcing used in this report.
\url{https://www.math.uni-bonn.de/people/schlicht/Prikry/a_minimal_prikry_forcing.pdf}.

\bibitem{Goldberg}
Gabriel Goldberg.
\emph{The uniqueness of elementary embeddings}.
Journal of Symbolic Logic \textbf{89} (2024), 1430--1454.
\doi{10.1017/jsl.2024.60}.
\arxiv{2103.13961}, version~2, Theorem~4.8 (Vop\v enka).
The theorem numbering here refers to the arXiv version.
\url{https://arxiv.org/html/2103.13961v2}.

\bibitem{BCGHP}
Tom Benhamou, James Cummings, Gabriel Goldberg, Yair Hayut,
and Alejandro Poveda.
\emph{Compactness phenomena in HOD and the optimality of Magidor's covering theorem}.
\arxiv{2608.24190}, version~1, 25 August 2026.
Theorem~1.1 gives the related optimal measurable upper bound for the first
power-set disagreement at $\aleph_\omega$.
\url{https://arxiv.org/html/2608.24190v1}.

\bibitem{ABGL}
Juan~P. Aguilera, Joan Bagaria, Gabriel Goldberg, and Philipp L\"ucke.
\emph{Large cardinals beyond HOD}.
\arxiv{2509.10254}, version~1, 2025.
Proposition~3.9 gives the low-rank HOD coding used in the dossier's
ultraexacting applications; Theorem~4.5 gives the published
$\Izero$--ultraexacting equiconsistency result.
\url{https://arxiv.org/html/2509.10254v1}.

\end{thebibliography}
\end{document}
